\documentclass[10pt]{article}
\usepackage[margin=20mm]{geometry}
\usepackage{amsmath,amssymb,fontspec,unicode-math,xurl,hyperref}
\IfFontExistsTF{DejaVu Serif}{\setmainfont{DejaVu Serif}}{\setmainfont{Latin Modern Roman}}
\setmathfont{latinmodern-math.otf}
\setlength{\parindent}{0pt}\setlength{\parskip}{6pt}\sloppy
\title{Technical Verification of Modified K-means Algorithm with Local Optimality Guarantees}
\author{AI-assisted technical review}
\date{}
\begin{document}
\maketitle
Abstract: The principal Bregman descent identity and the finite-convergence argument are supported by analytic inspection and independent small checks. The second assertion of Theorem 4.3 needs its intended partial-optimality premise. A pinned later implementation omits a crucial empty-cluster repair condition and admits a nonlocal C-LO termination on a three-point example. The theoretical algorithm excludes that trace. The supported assessment is Minor, with the published experiments’ implementation provenance unresolved.\par
\section*{Identity and scope}
Mingyi Li, Michael R. Metel and Akiko Takeda, Modified K-means Algorithm with Local Optimality Guarantees, ICML 2025, PMLR 267:35656–35684, I25-43694 (0325), OpenReview `v68wPjgEQ8`. The frozen publisher PDF has 29 pages; all page references use its printed/PDF numbering. The main text and Appendices A–F were read. The bibliography was not independently audited. Figures 1–3, Table 1 and Tables 10–11, together with decisive theorem/proof pages, were visually inspected.\par
The available repository is pinned at commit `7493c390670da6071b2d56bb937a4812a3ef8549`, dated 13 February 2026. The complete Python and C++ algorithm files, README, C++ settings/statistics helper and real-world experiment driver were read statically. The two other run entry points were not read. No author program, dataset experiment or benchmark was run. Separate OpenReview attachments were inaccessible; later errata were not surveyed. This later repository is not proof of which program produced the 2025 tables.\par
\section*{Definitions and supported derivations}
The data are distinct points in the interior of the convex domain of a strictly convex differentiable generator \(\phi\), with positive weights, \(N>K\), and Bregman divergence \(D_\phi(x,c)\). The hard assignment set \(S_1\) consists of vertices of the product-simplex relaxation \(S_2\). For \(F(P)=\min_C f(P,C)\), nonempty clusters have unique weighted-mean centers. \(A(C)\) is the set of minimizing assignments for fixed centers, while \(T(P)\) contains the hard assignments obtained by changing one point’s cluster. D-local optimality compares \(F(P)\) with all of \(T(P)\); C-local optimality is local optimality of \(F\) on \(S_2\). These are different from global optimality.\par
The concavity of \(F\), as an infimum of linear functions of \(P\), supports the D-local to C-local implication: near a vertex, feasible directions lie in the cone generated by its one-point neighbors, and concavity bounds the objective below by convex combinations of their values and \(F(P)\). The nonempty-cluster condition is essential when using continuous dependence of the weighted means on \(P\).\par
Lemma 4.1 correctly expresses the change from moving an \(\alpha\) fraction of point \(g\) from \(a\) to \(b\): \[\begin{aligned}
 F(P')-F(P)={}&\alpha w_g\{D(x_g,c_b)-D(x_g,c_a)\}\\
 &-(s_a-\alpha w_g)D(c'_a,c_a)-(s_b+\alpha w_g)D(c'_b,c_b).\end{aligned}\] The weighted-mean Bregman decomposition proves this identity. Independent numerical evaluations with squared, KL and Itakura–Saito generators on positive scalar data agree with direct recomputation to within \(10^{-12}\). This finite check supplements, rather than replaces, the algebra. Inspection of the original p. 16 confirms that Lemma D.2 has the correct \(p'_{k,n}-p_{k,n}\) sign; an initially suspected extraction-related sign problem is rejected.\par
The paper’s example \(X=(-4,-2,0,1.5,2.5)\), unit weights and \(K=2\) also checks. At centers \((-2,2)\) and the indicated \(3+2\) partition, \[F(\alpha)=\frac{57}{2}-\frac{36}{3-\alpha}-\frac{16}{2+\alpha},\qquad
 F'(\alpha)=\frac{-20\alpha(\alpha+12)}{(\alpha-3)^2(\alpha+2)^2}<0
 \quad(0<\alpha\leq1).\] Thus \(F(0)=17/2\) and \(F(1)=31/6\), validating this failure of ordinary K-means to reach a C-local solution under the paper’s definition.\par
For Theorem 4.2, a partial optimum with all masses positive and a unique nearest assignment has strictly positive assignment gaps. They persist for centers in a neighborhood. The weighted means \(C(P')\) are continuous there, so \(f(P',C(P'))\geq f(P,C(P'))\geq F(P)\) for nearby \(P'\). This independently supports C-local optimality. Lemma D.4 establishes the partial-optimality and distinct-center premises at the algorithm’s special branch; strict descent through a finite set of hard assignments then supports Theorem 4.4’s finite termination. The complexity result concerns each iteration, not an equal total runtime or a polynomial iteration bound.\par
\section*{Findings}
Problem 1 (Theorem 4.3’s second assertion lacks a premise).  The statement on p. 5 and its proof on p. 18 assert that nonsingleton \(A(C_P^*)\) implies \(A(C_P^*)\cap T(P)\ne\varnothing\), assuming only distinct centers for a solution \((P,C_P^*)\). Partial optimality is also necessary for this assertion as written. This is a Minor, local statement defect.\par
Take \(\phi(x)=x^2\), six distinct real points \((-5,0,2,-2,1,4)\), unit positive weights and \(K=2\). Assign \(\{-5,0,2\}\) to the first cluster and \(\{-2,1,4\}\) to the second. Both clusters are nonempty, their unique optimal centers are \(-1\) and \(1\), and \(F(P)=44\). The point \(0\) ties between these centers, so \(A(C_P^*)\) has two elements. However, point \(2\) is assigned to the wrong nearest center, as is point \(-2\). Every member of \(A(C_P^*)\) must change both assignments, and hence differs from \(P\) in at least two coordinates. The intersection with \(T(P)\) is empty. Exact enumeration of the \(2^6\) assignments gives Hamming distances \(2\) and \(3\) for the two minimizers.\par
Adding \(P\in A(C_P^*)\), equivalently partial optimality here, repairs the second assertion. A tied point can then be moved alone while every other assignment stays minimizing. Distinct centers ensure that this tied point differs from each center, and Lemma 4.1 gives strict decrease through the positive center-shift terms. The first assertion, that a C-local solution with distinct centers has singleton \(A(C_P^*)\), is not refuted by this example. The text immediately below the theorem discusses converged K-means, which supplies the missing premise. Lemma D.4 likewise supplies it at the actual invocation in Theorem 4.4. Thus this counterexample does not invalidate the main algorithm’s convergence proof. Reading the word “solution” as implicitly meaning “partial optimal solution” would also avoid the defect, but that qualification is explicit in Theorem 4.2 and absent in 4.3; stating it removes the ambiguity.\par
Problem 2 (Released C-LO can terminate outside the proved conditions).  The pinned Python and C++ implementations omit Algorithm 1’s \(x_g\ne c_b\) requirement when repairing an empty cluster. Their random initialization samples with replacement. Together with the tie handler, this admits non-C-local output. This is a confirmed release-level defect; its historical experimental effect is Needs Evidence.\par
For the C++ file `cpp/LO\_K\_means.hpp`, random initialization at lines 81–84 draws independently, empty-cluster repair at 222–229 checks only that the donor has more than one point, and lines 242–244 assign a tied point to the second minimizing index without ensuring a change. The corresponding Python paths are at 79–83, 165–170 and 173–180. The public API permits ordered distinct data \(X=(0,-1,1)\), unit weights, \(K=2\), squared divergence and random initialization. Both sampled centers equal \(0\) with probability \(1/9\) under independent uniform draws.\par
At these centers all points tie, so nearest assignment sends all three to cluster 0. The empty-cluster repair takes the first point, \(0\), into cluster 1. Both resulting centers remain \(0\) and the labels are \((1,0,0)\). The next iteration repeats those assignments and repair. The tie handler examines point \(0\), sets its label to 1, which is already its label, and breaks. The unchanged-assignment stopping condition terminates with loss \(2\). This is a static state trace of the inspected code, not an execution of the authors’ program.\par
Move an \(\alpha\) fraction of point \(-1\) from cluster 0 to cluster 1. The cluster masses are \(2-\alpha\) and \(1+\alpha\), both positive, with centers \(\alpha/(2-\alpha)\) and \(-\alpha/(1+\alpha)\). Exact recomputation gives \[F(\alpha)=2-\frac{\alpha^2}{2-\alpha}-\frac{\alpha^2}{1+\alpha}<2
 \qquad(0<\alpha\leq1).\] Assignments and centers converge to the returned state as \(\alpha\downarrow0\); it is therefore neither C-local nor CJ-local. At \(\alpha=1\) the loss is \(1/2\), so it is not D-local either. The supplied input satisfies the data, weight and divergence assumptions; the failure comes from implementation choices that differ from the paper’s algorithm.\par
The paper samples without replacement and explicitly requires \(x_g\ne c_b\) in the empty-cluster repair. Its proof relies on the resulting strict decrease. These conditions exclude this trace and are material counterevidence against treating it as a false theoretical guarantee. The real-world driver selects random initialization, but its helper sorts and deduplicates data; that driver does not preserve the particular input ordering above. No assertion is made that a reported benchmark encountered this state. Aligning initialization and repair with the pseudocode, and requiring a genuine tie move, addresses the exhibited mechanism. A fully verified finite-precision implementation and historical rerun remain open. The shared \(10^{-10}\) tolerances and positive-mass thresholds also mean the release is not an exact-arithmetic implementation for arbitrarily small positive weights.\par
\section*{Impact and empirical limits}
These findings leave the principal descent/convergence argument intact under its intended partial-optimality premises and exact pseudocode. The later release requires repair before its output can be given the same unqualified guarantee. Figures 1–3 and the real-data tables give descriptive comparisons; raw runs and their uncertainty were not reproduced. The released implementation stores an \(N\times K\) distance array, whereas the paper says an \(O((N+K)d)\)-space implementation is possible. This release choice does not disprove that existence claim.\par
Dimension | State\par
Overall / theorem qualification | Minor / Minor\par
Released C-LO counterexample | Confirmed by static trace and exact objective\par
Repair scope | Theorem premise supported; implementation rerun open\par
Published experimental impact | Needs Evidence\par
Adjudication / author adoption | Pending / unknown\par
Human review | None recorded\par
2 M. Li, M. R. Metel and A. Takeda, Modified K-means Algorithm with Local Optimality Guarantees, ICML 2025. https://proceedings.mlr.press/v267/li25bt.html. Authors’ repository, commit `7493c390670da6071b2d56bb937a4812a3ef8549`. https://github.com/lmingyi/LO-K-means.\par
\section*{Evidence record}
Source SHA-256: `816b67b19ae88db2abdfd5490aec191a7f04fc691f07c8443b0686023c4ab1f4`. The accompanying independent Python 3.13.9 script uses rational arithmetic, SymPy 1.14.0 and NumPy 2.4.3 for the specified enumeration, symbolic losses and three Bregman identity checks. It imports no author code. These are initial self-checks, not independent acceptance or human review.\par
\end{document}
